\documentclass[11pt]{article}

\def\chapitre{d'introduction.}
\def\pagetitle{Individu et Communauté.}

\input{setup.tex}

\begin{document}

\input{title.tex}

\thispagestyle{fancy}

\begin{center}
    judith.ertel2@gmail.com
\end{center}

Banwords : De plus; Dans un premier temps...; D'abord, Ensuite, Enfin, on retrouve.

\section{Introduction.}

Les termes << Individu >> et << Communauté >> sont manifestement en opposition :\\
--- Du point de vue moral, est-ce que l'individu prime sur la communauté ou l'inverse ?\\
--- Du point de vue juridique, est-ce que l'individu est façonné par la communauté ?\\
$\hra$ Jusqu'aux années 70, l'homosexualité était punie par la communauté.

Qu'est-ce qu'un individu ?\\
$\hra$ Un individu est ce qui ne peut être divisé. Ce terme naît tardivement, quand le soi devient primordial, vers le XVIIe siècle, notamment avec Descartes. Ce terme engage de nombreuses questions : la place que l'on a au sein de la société, nos libertés. Avoir constamment conscience que la communauté humaine existe est un trait propre à l'humain.\\
$\hra$ Ce qui singularise l'individu, c'est l'ensemble des caractéristiques qui lui sont propres, par sa manière d'être. C'est ce qui est identifiable à lui seul, à l'exclusion de tout autre.\\
$\hra$ L'\bf{individualisme} est un mode de pensée plaçant les intérêts personnels au-dessus de ceux de la communauté.\\
$\hra$ L'\bf{individualité} est la distinction entre les individus.\\
$\hra$ Le \bf{holisme} : l'individu ne peut être considéré que dans sa globalité.

Or cet individu n'existe pas sans faire parti d'un tout. Est-ce que les individus existent bien par eux-mêmes, en singularité ? Que devient l'individu en dehors de toute communauté ? L'individu n'existe-t-il que pour défendre, et appartenir à une communauté ?

Qu'est-ce qu'une communauté ?\\
$\hra$ C'est un ensemble plus ou moins large, qui intègre les individus, et ces individus ont donc en commun leur univers symbolique, leurs valeurs, leurs codes, et une relation affective. Cela suppose qu'il y ait du commun entre les individus, ce qui est différent de la société. 

Comment peut-on se constituer à la fois en étant au sein d'une communauté et à l'écart ? Il faudrait que la communauté ne soit pas excessivement restrictive.
Les communautés principales sont la communauté familiale, ethnique, religieuse, politique.

Aristote dit que la cité est l'organisation rassemblant les hommes en vue du bien commun. L'Homme est alors forcément un animal politique, il est d'abord un être civique, il est citoyen de l'organisation. Il cherche à démontrer que la cité est antérieure à l'Homme : il n'y a pas d'Homme en dehors de la cité. Il dit que c'est une necessité naturelle, si un homme est seul, il est soit supérieur : un dieu, soit dégradé. L'Homme est intégré à la cité, reposant sur le logos, pour développer des valeurs communes. 

Les communautés se forment aussi autour de récits fondateurs, de mythes, par exemple dans l'Empire Romain, autour de la fondation de Rome par Romulus et Rémus.\\
$\hra$ Un mythe est un récit fictif ou qu'on ne peut pas prouver, qui a une fonction symbolique.

Bouc-émissaire : on fait porter à une personne, un groupe, une communauté la responsabilité des maux qui atteignent une communauté, qui cherche une victime expiatoire pour dépasser une crise, ressouder la communauté.\\
$\hra$ Les juifs ont été longtemps bouc-émissaire, notamment lors de la seconde guerre mondiale.\\
$\hra$ Les femmes ont été longtemps bouc-émissaire, notamment lors de la chasse aux sorcières.

Comment l'individu peut-il développer sa subjectivité si la puissance psychique de la communauté le gagne ?\\
Les rites, les histoires, les préceptes lui sont transmis depuis l'enfance, malgré lui, donc il est ce que l'on a forgé de lui. Il a toujours vécu dans cette communauté, ces règles deviennent donc l'intimité de l'individu, et donnent l'illusion que c'est sa subjectivité, il ne se rend pas compte que cette manière de voir le monde provient de l'extérieur, du regard de l'autre, du jugement moral, de la fidélité, qui sont des éléments extérieurs. C'est la pression de la communauté qui pousse à se conformer aux règles.

À partir du XVIIIe siècle, l'individu se pense autrement, d'abord avec Descartes : << Je pense donc je suis >>, l'idée que l'individu a une existence singulière parce qu'il a une conscience réflexive émerge, alors qu'il se pensait jusque là comme membre d'une société, d'une famille... Avec l'idée que ce qui compte, c'est le bien-être des individus, la liberté de penser. Cela vient des guerres de religion, après la naissance du protestantisme. L'idée de tolérance apparaît. Les gens qui ont défendu la tolérance, la liberté ont fait que les individus ne soient plus sous un pouvoir (religieux, ...).\\
$\hra$ Affaire Calas : un individu ne peut pas être responsable de sa famille.

Rousseau, Hobbes et Locke ont pensé le rapport entre individu et communauté différemment, Rousseau pensait que l'Homme isolé était bon, et que les relations sociales l'ont corrompu. Pour Hobbes, << L'Homme est un loup pour l'Homme >>, le contrat social serait là pour éviter la violence, il faut que la société s'organise pour réaliser un bien être économique avec le produit du travail des individus, qui doit être encadré, pour protéger la propriété privée. L'individu devient alors celui qui défend le produit de son travail, sa propriété privée. On demande à la société politique de nous apporter le maximum de bien-être, l'individu demande ça à la société, et l'individu apporte à la société son travail.
L'individu veut que la société lui permette à lui seul d'être bien, et lui demande de plus en plus de choses. Elle doit améliorer la condition des travailleurs, des femmes, améliorer les loisirs... C'est l'État providence. L'individu est universel, il peut être émancipé de tout groupe sur le plan juridique, il n'est pas tenu d'avoir une religion, on ne peut pas le considérer au regard de sa religion (en France) 

Lorsque l'individu se détache de toute communauté, il y a deux risques : l'isolement, la société contemporaine a vu se renouveler la question de la communauté et des communautés religieuses.\\
L'universalité est de considérer que l'individu est irréductiblement que dans une seule communauté qui surplombe toutes les autres : la communauté humaine. 


\section{Eschyle.}

On parle d'une période archaïque jusqu'au 6e siècle avant J.C., c'est une période de mise en place de systèmes politiques, appelée tyrannie, qui a fait que le territoire de la Grèce Antique allait de la Mer Noire. Il y a eu la constitution des cités à ce moment, et la tyrannie s'est vue écartée du pouvoir par des communautés de citoyens-soldats. C'est une sorte d'aristocratie militaire qui tenait ces territoires, ce qui a fondé la cité grecque, qui a évoluée au fil des siècles. C'est à partir de ce moment que les \bf{citoyens} sont sur un pied d'égalité, on parle de période Classique. Il y a alors plusieurs systèmes politiques : l'oligarchie, tenue par un petit groupe de personnes; et la démocratie grecque, créée par Clisthène (c. -570), qui créée des institutions dans lesquels les citoyens se représenteront, avec une magistrature, un corps électoral et un corps politique éxecutif. La cité a une assemblée du peuple qui se réunit à intervalles réguliers, qui est souveraine, qui désigne les magistrats. Il y a un conseil, la Boulè, constituée de citoyens tirés au sort, décidant d'éléments politiques : projets de lois.
Dans cette période, il y a eu les guerres médiques en Perse, il y a eu notamment la bataille de Marathon et la bataille navale de salamine. En sortant victorieux de ces guerres, les grecs ont eu un élan religieux, et ont reconstruits des temples, ... Cela a aussi permi un essor du théâtre et de la littérature, où les acteurs représentaient les héros de guerre, d'autant plus pendant les fêtes religieuses, et les grandes dyonisies. Il y avait quatre types de théâtre : le dithyrambe, de grands poèmes en l'honneur d'un Dieu ou d'un héros, chantés par un choeur; la tragédie; le drame satyrique, une sorte de sujet mythique, plus léger que la tragédie; et la comédie, faite par un cortège. Le premier concours de tragédie se serait tenu en -534 pendant les grandes dyonisies. Il ne reste que trois noms d'auteurs: Eschyle, Sophocle et Euripide; et que 33 pièces restantes, sur plus de 1500.

Dans le théâtre, il y a le Choeur, représentant le peuple, commentant les actions, mais n'agissant jamais dans la pièce. Les personnages font avancer l'action, mais le Choeur jamais. Le Coryphée est le chef du Choeur. Le théâtre est composé d'un scène, d'une proscène, de l'orchestra, et de gradins pour le public. Les acteurs étaient dans des costumes très lourds avec des masques. C'est souvent les frottements entre ce que disent les acteurs et ce que dit le Choeur qui mettent en relief l'individu et la communauté.

On a affaire à une confrontation entre les femmes et les hommes, entre ces deux communautés, deux individus hommes : Danaos et Pelasgos, qui sont en fait auprès de la communauté des femmes, qui n'est pas individualisée, qui forme un tout homogène et unique. Dans les Sept contre Thèbes, le Choeur est composé de femmes, et la guerre est une affaire d'hommes. Il y a cette interrogation sur l'alterité à travers cette question des hommes et des femmes, du pouvoir, de l'étranger et de l'ennemi.
Trois questions conditionnent les communautés : la question ethnique, la question sexuelle et raciale, mais aussi l'individu, car c'est le langage qui va conditionner l'individu, qui lui permet de l'affilier à une communauté. Les femmes se présentent comme des victimes de la menace des fils d'Egyptos.
Le roi constate des Suppliantes ayant en commun des rites, des croyances et des symboles communs, il est alors d'accord pour accepter les argiennes. Pelasgos accède à la parole des femmes et leur demande de lui faire confiance et d'avoir une parole libre.

\bf{Les Suppliantes.}

Il reste d'Eschyle, qui avait écrit près de 90 piècess, uniquement 7 pièces. Les Sept Contre Thèbes sont datées de -467, et Les Suppliantes de -464. Il a gagné 13 concours aux Dyonisies.

Zeus s'éprend d'une argienne nomée Io, prêtresse d'Héra. Or Héra était amoureuse de Zeus et était furieuse d'être ainsi trahie. Elle transforme alors Io en génisse. Zeus se transforme alors en taureau. Héra lance alors un taon contre la jeune génisse, et fuit pour l'éviter. Elle franchit le détroit du Bosphore, et va jusqu'en Égypte, où Zeus va la sauver, et la retransformer en humaine, et met au monde Epaphos, créant une lignée, parmi ses fils : Danaos qui a 50 filles et Egyptos, 50 fils. Egyptos voulait marier ses fils aux filles de Danaos pour obtenir leurs dotes. Danaos refuse car une oracle lui avait dit qu'il se ferait tuer par un de ses gendres. Il pousse ses filles à aller à Argos pour chercher asile. Les Suppliantes s'ouvre au moment où elles arrivent au portes d'Argos. Elles paviennent à persuader Pélagos, le roi d'Argos, de plaider leur cause devant le peuple. Il parvient à persuader son peuple, en leur cachant des dangers de la guerre.

Il perd la guerre, et Danaos est contraint de marier ses filles à ses neuveux. Il leur ordonne de tuer leurs maris, et toutes lui obéissent, sauf une : Hypermestre, qui épargne son mari, Lynceus, qui tuera Danaos. Les 49 autres filles sont condamnées à porter un vase percé pour éternellement remplir un tonneau percé.

\bf{Analyse.}

Le Choeur est constitué des Danaïdes, elles se mettent sous l'autorité de Zeus et appellent aux valeurs de la cité : la justice et le droit. Elles s'opposent d'emblée à ce que les hommes les contraignent. On a la communauté des Danaïdes et des hommes qui s'opposent.

Danaos prévient ses filles des conditions de l'hospitalité : la prudence, le respect des Dieux, les raisons de l'exil, ...

Immédiatement, sur leur apparence, les Danaïdes sont vues comme étrangères : <<basané>>. Or la parole permet de faire le lien avec les étrangers : << tu as la parole >>. Il s'agit par la parole de voir que leurs origines sont communes. L'apparence n'est donc pas celle des origines. Donc les hommes, en faisant foi à la parole, peuvent dépasser les apparences.

Ce debut de pièce pose un discours critique de la communauté. 

\bf{Les Sept contre Thèbes.}

Cadmos est le fondateur de Thèbes, il a un fils Laïos. Pelops, roi de Pise, confie son fils Chrysippe à Laïos. Il se rend compte qu'il est amoureux de lui, et le viole. Chrysippe se suicide, et Pelops demande à Apollon de maudire Laïos.\\
Laïos se fait tuer par son fils Oedipe, et Oedipe épouse Jocaste. Ils ont deux fils et deux filles. Des années après, il se rend compte qu'il a tué son père et épousé sa mère, il se crève les yeux, ses fils refusent de prendre soin de lui, et Oedipe les maudit en disant qu'ils mourront en s'entretuant.\\
Étéocle et Polynice essaient de conjurer la malédiction en se mettant d'accord sur un pouvoir tournant : ils règnent à tour de rôle sur Thèbes. Quand arrive le tour de Polynice, Étéocle refuse de lui donner le pouvoir, et Polynice fait appel à la cité d'Argos pour le défendre. Thèbes va l'emporter, mais les deux frères s'entretuent. Pour rétablir la paix, le frère de Jocaste, Créon prend le pouvoir et décide d'honorer Étéocle, et de laisser le corps de Polynice sans sépulture. Il fait sortir Polynice de l'humanité.
Une des filles, Antigone, s'oppose à la décision de son oncle, et se révolte, elle tente d'ensevelir son frère. Créon l'enmure pour la punir, et la malédiction se termine.

\bf{Analyse.}

La pièce commence sur Étéocle, individu, s'adressant à la communauté. Il en appelle au peuple de Cadmos. La pièce s'ouvre sur un siège, une menace pour toutes les institutions de la communauté. Pour protéger le peuple, il faut protéger les Dieux, les autels. La religion est une institution qui fait communauté, et il appelle à l'engagement de chacun : << chacun enfin se donnant au rôle qui convient à ses forces >>. Il dit lui-même que si ils gagnent, la communauté est préservée, tout comme la religion, car la victoire serait grâce aux dieux, la défaite serait de sa faute; sur le principe du bouc-émissaire, pour soulager la communauté. Il demande à tous de s'engager, mais c'est lui qui est à la manoeuvre. La communauté aurait donc besoin d'individus pour exister, sous la forme de guide suprême, de leader, ou de bouc-émissaire. À cela va s'ajouter une problématique qui est aussi dans Les Suppliantes : le Choeur, car quand il y a la guerre, c'est la guerre des hommes. Et le Choeur dans Les Sept Contre Thèbes est composé de femmes, qui refuse la légitimité de la guerre.



\pagebreak

\section{Spinoza.}

\bf{Cours d'introduction, Spinoza.}

Opposition entre individu et communauté retrouvée dans la philosophie du XVIIe / XVIIIe siècle, Spinoza est un philosophe comme un autre à parler de la dialectique entre individu et communauté. Au XVIIe siècle, l'idée qu'il y a des individus hors-communauté, à l'état de nature, est communément pensée. Aujourd'hui, un individu à l'état de nature, on pourrait dire qu'il n'y en a plus, et même qu'il n'y en a jamais eu, ou bien c'est une illusion, ou une fiction théorique vis-à-vis de laquelle il ne faut pas donner de prétention historique. La figure de l'individu à l'état de nature est donc illusoire.

La notion d'individu renvoie à l'idée d'une division, il y a une totalité de laquelle l'individu participe, est la partie divisée. Un individu est la dimension irréductible du point de vue de la totalité, il n'y a pas d'individu sans totalité. L'individu est une personne humaine singulière.
On peut considérer qu'elle a elle-même une nature / une essence propre, pas conditionnée d'un ensemble social auquel cette personne humaine participe. Il y a une réalité personnelle, qui existe indépendemment de la société, une nature humaine préalable à l'être social. C'est un axe sur lequel Spinoza met l'accent, avec le droit naturel, l'Homme a des lois naturelles qui lui sont propres. Une bonne société, politiquement bien constituée, préserve le droit naturel de chacun.
L'autre façon de considérer l'individu est qu'il est essentiellement conditionné, formé par la réalité sociale dans laquelle elle se trouve. Aristote présente l'Homme essentiellement social, ou politique, ne peut pas ne pas être politique. Aristote dit qu'un Homme qui ne serait pas politique, est ou un Dieu, ou une bête. 

Spinoza dit Dieu ou la nature, << Deus sive natura >>, pour lui Dieu est la nature, Dieu n'est pas la religion, une religion est une institution humaine. Il se demande comment Dieu dicterait une bonne conception politique. Cela renvoie à l'épisode de Moïse. Toute la question des lois naturelles se trouve comprise dans cette conception de Dieu. Spinoza veut dire qu'il y a une nécessité à l'oeuvre dans la nature qui n'admet aucun arbitraire, aucune volonté capricieuse qui permettrait de déroger à ces lois. Dieu n'admet aucun caprice ou miracle, la conception du Dieu de Spinoza est opposée à celle d'un dieu chrétien, ce dieu peut pas vouloir autrement que ses lois. La réalité ne peut pas s'affranchir de ses lois, même par les miracles. Le Dieu chrétien reste toujours mystérieux à l'intelligence humaine, elle est limitée alors que Dieu est infini. Descartes dit que l'entendement humain est infini dans l'absolu, mais la difficulté est que la vie de chacun est limitée. Dans ce cadre là, la compréhension de l'infinité de Dieu est impossible. Il est donc impossible de connaître la totalité de ce qui régit la nature : << Les voies du seigneur sont impénétrables >>. Les miracles sont sensées être des preuves du pouvoir de Dieu sur la nature. Toute puissance théologique, théocratique prend appui sur cette crédulité humaine, elle a pour principe de diminuer la puissance rationnelle de l'Homme au profit de sa puissance de crédulité. Spinoza décrit que lorsque le sort est favorable, l'homme est conscient dans sa seule raison, dans la dimension rationnelle du monde, en revanche, quand il est défavorable, quand l'homme ne trouve pas de solution, il renonce a sa raison, et verse dans la superstition, à la possiblité d'une intervention divine, malgré la rigueur des lois de la nature.

Dans une République, Spinoza dit qu'on ne doit pas faire économie de la religion, qui a une fonction d'adhésion, or on a l'impression que plus le temps avance, moins on fait référence à la religion, alors que pour Rousseau, et Spinoza, c'est un instrument très efficace à la cohésion des individus dans la communauté, dès lors qu'elle respecte sa juste place.  

Ces lois de la nature, comment Dieu nous les fait-il connaître ? D'abord, par la science, la connaissance scientifique de la nature, est la connaissance de Dieu. C'est une idée très présente dans le courant rationnaliste. C'est l'idée que la nature est un grand livre écrit en langue mathématique, et pour apprendre à la connaître, il faut apprendre cette langue mathématique. Aujourd'hui, Dieu est opposé à la connaissance scientifique, c'est une contradiction très violente. En quoi ces lois de la nature permettent-elles d'en savoir plus sur les lois humaines ? Spinoza emprunte l'exemple de la République des Hébreux, pour montrer à quel point le conflit entre l'autorité politique et religieuse a été destructeur pour l'État lui-même. Il propose une analyse de ce conflit. Il est illustré par la figure de Moïse, importante du point de vue religieux et politique. Or, cette autorité de Moïse lui est confiée par Dieu, qui fait connaître ses lois par les dix commandements. Moïse fait état de son rapport immédiat à Dieu, il devient son médiateur auprès des Hommes, qui vont avoir du mal à reconnaître l'authenticité de ce rapport immédiat à Dieu. L'Homme est un être de nature, comme tous les êtres, son existence n'est en rien contradiction avec les lois de la nature en général. 

Tout être par nature, désire persévérer dans son être, et d'une certaine manière, il persévere dans son être aussi longtemps que la puissance d'un autre être qui désire lui aussi persévérer dans son être, ne contredit pas ce mouvement, ça c'est le principe de tous les êtres dans la nature. L'Homme possède des caractères ou des puissances qui lui sont propres, et en particulier, dans sa nature, dans son essence, il y a un caractère, une dimension qui lui est propre, non pas les émotions, pour Spinoza, il n'est pas exclu que les animaux aient des émotions, c'est la raison, la capacité de raisonner, et de vivre rationnellement, sous les lois de la raison. Pour Spinoza, cette vie rationnelle est tributaire d'une sagesse, d'une morale, d'une rationnalité. 

Pour Hobbes, c'est la peur d'être confronté à autrui, et que l'autre le tue, qui est au fondement du projet social, justifiant << Tu ne tueras point >>, lui donnant sa légitimité en tant que loi divine, tout comme les autres commandements. La société a pour projet d'interdire et de conjurer cette peur de tous contre tous. Hobbes a la conviction que la société humaine est construite en réaction à une situation dans l'état de nature où les hommes s'entretuent. Pour Spinoza, il y a une rupture dans l'Histoire entre l'état de nature, et la société instituée, pour Hobbes, c'est seulement à partir du renoncement par chacun à sa puissance personnelle. La société pour Spinoza serait dans la continuité du droit naturel de chacun, tout en apportant des solutions aux problèmes des Hommes.

\bf{Cours.}

Quelle place est laissée par l'État à l'individu, et quel régime politique est le meilleur pour l'individu ? L'individu est gouverné par ses affects par nature : plaisirs personnels. Une société va apporter un cadre juridique sur cela. Si l'Homme a un mécanisme qui relève de son droit naturel avec l'idée qu'il ne peut pas être autre chose que lui-même. Tout être, d'après Spinoza, persévère dans son être, il ne peut pas être autre chose. Aucun Homme ne peut détourner l'être de ce qu'il est, sauf s'il pense que quelque chose lui est plus utile. L'Homme fera tout pour ce qui lui fait du bien. L'homme seul ne peut rien construire, il comprend par la raison que ce qu'il y a de mieux pour lui est de se mettre en société avec d'autres, avec un pacte social, c'est une organisation politique qui lui enlève deux choses : le sentiment qu'il pourrait se détruire et les sentiments liés à la crainte. Car sinon, il s'en remet à Dieu; or pour Spinoza, Dieu est la nature, donc il ne sert à rien de s'en remettre à Dieu. Pour Spinoza, la démocratie assure la paix, les autres systèmes politiques nient aux individus la possibilité d'être libre, et créent des instabilités, donc l'Homme va spontanément aller contre l'absolutisme, et pousser l'État à agir sans empiéter sur sa liberté d'individu parce qu'il sait ce qui est le plus utile pour lui-même. Dans un état absolutiste, le pouvoir impose des choses à l'individu, dans la démocratie, les individus transfèrent une part de leur pouvoir à l'état. L'intérêt commun aide l'intérêt privé. Quand l'État est stable et aide les individus à développer leur puissance, les individus n'ont pas besoin d'agir selon des affects et appétits non rationnels. Si l'Homme est pris uniquement par ses appétits, il se trompe sur ce qui est bon pour lui. 
Un citoyen est un sujet qui a consenti rationnellement à consentir aux lois de l'État. Spinoza pense que les lois ne sont pas forcément morales, mais sont là pour nous préserver.
Le pacte avec Dieu vient d'une révélation, il n'est pas naturel. Il faut qu'il y ait un intermédiaire, une révélation, aucun raisonnement ne peut établir l'existence de Dieu. Il faut des signes. Il consiste à transférer ses droits à Dieu : renoncer à ses droits naturels.
La raison est contrée par des forces : les désirs, les appétits, les passions. La seule résistance est la notion de consentement à la délégations des droits. Il faut que l'individu consente à déléguer ses droits, et à accepter que tout n'est pas fait pour soi individuellement.

\bf{...}

La sphère privée semble permettre à l'individu de se libérer de la pression sociale, apportant un cadre de liberté ou d'émancipation, mais cela peut être une illusion.
Dans Les Sept Contre Thèbes, Antigone rejette la sphère politique, elle considère que les deux frères ont le même droit à la sépulture, relevant de la sphère privée. 
Spinoza ne s'intéresse pas à la sphère, mais développe l'idée que la sphère privée est l'espace de la liberté religieuse. 
La morale ne doit pas être un des fondements politiques, la religion ne doit pas interferer avec la politique.

Dans l'oeuvre de Spinoza, l'Homme est présenté comme étant libre d'être lui-même mais sans le cadre politique, seul avec lui-même, il ne peut pas être libre d'exercer ses potentialités. Dans les textes d'Eschyle, les personnages combattent les injonctions de la sphère privée, notamment les Danaïdes qui ne peuvent se résoudre.

\section{Edith Wharton.}

Conventions : règles sociales, morale; convenance : ce qui est bien vu; conformisme : se prêter au groupe. Cela apporte une uniformisation : mêmes pensées, mêmes désirs : société normative.
Bon goût : distinction entre les classes, signe exterieur de comportement.\\
La mère de Archer pense que toute la structuration en cercles maintient une société fonctionnelle.

L'Opéra du début de l'histoire est comme un miroir des personnages de l'histoire. Ce groupe est une communauté sociale où tout le monde se connaît : il tient par une proximité culturelle, et affective. Ils sont unis par des valeurs culturelles et puritaines, selon des lois religieuses, sociales, fondées sur les institutions (le couple, la famille, la haute société). Comment l'individu trouve-il sa place dans la société ? Est-ce que pour exister vraiment, la seule solution est de s'échapper de la société ? Les personnages ont des habitudes fortes : il y a ce qui se fait, et ce qui ne se fait pas.\\
Cette société est hiérarchisée avec une organisation de plusieurs entités. Les personnages savant qu'ils appartiennent à un certain cercle, ayant des privilèges par rapport à d'autres, et moins que certains autres.

Newland Archer : adapté à cette société, le roman est la quête de comment concilier les injonctions du groupe et son individualité.
Ellen Olenska : origines new-yorkaise, a vécu en Europe; différence de marche entre les moeurs new-yorkaises et européennes. Elle n'adhère pas aux conventions New-Yorkaises : divorce. Ellen choisit d'habiter dans un quartier hors de la communauté : << son petit refuge >>.



\end{document}